\section{Описание}

В третьей лабораторной работе исследовалось качество программы-словаря из лабораторной работы \textnumero 2. Для этого были выделены две версии решения:
\begin{itemize}
\item \texttt{baseline} --- точная копия программы из предыдущей работы;
\item \texttt{fixed} --- версия после исправления найденного недочёта.
\end{itemize}

Основная цель состояла не в повторной проверке корректности красно-чёрного дерева как такового, а в исследовании практических свойств готовой программы: времени выполнения операций и поведения при работе с динамической памятью.

\subsection{Подготовленные средства}

Для исследования были подготовлены следующие утилиты и вспомогательные программы:
\begin{itemize}
\item \texttt{profile.cpp} --- изолированные замеры операций \texttt{insert}, \texttt{find}, \texttt{erase}, \texttt{save}, \texttt{load};
\item \texttt{benchmark.cpp} --- сравнительный workload-бенчмарк в стиле лабораторной работы \textnumero 2;
\item \texttt{gen\_profile\_data.py} --- генерация детерминированного набора из уникальных слов и значений;
\item \texttt{run\_quality\_suite.sh} --- запуск серии измерений с записью результатов в \texttt{CSV};
\item системные утилиты macOS \texttt{sample}, \texttt{leaks}, \texttt{heap}.
\end{itemize}

Для профильных запусков использовался набор из \(200000\) уникальных ключей. Для \texttt{insert}, \texttt{find} и \texttt{erase} измерялись только операции над деревом; подготовка словаря для режимов \texttt{find}, \texttt{erase}, \texttt{save} и \texttt{load} выполнялась вне измеряемого участка. Временные значения в таблицах ниже приведены как медиана по пяти запускам.

\subsection{Дневник выполнения}

\begin{longtable}{|p{0.17\textwidth}|p{0.77\textwidth}|}
\hline
Дата & Выполненные действия \\
\hline
04.04.2026 &
Создана отдельная директория \texttt{lab3}. В неё перенесена базовая версия словаря из лабораторной работы \textnumero 2 и выделены две ветви исследования: \texttt{baseline/main.cpp} и \texttt{fixed/main.cpp}. \\
\hline

04.04.2026 &
Подготовлены программы \texttt{profile.cpp}, \texttt{benchmark.cpp}, генератор \texttt{gen\_profile\_data.py} и сценарий \texttt{run\_quality\_suite.sh}, чтобы отделить исследовательскую часть от исходного решения. \\
\hline

04.04.2026 &
Сгенерирован детерминированный набор данных из \(200000\) слов. На нём получены baseline-замеры для \texttt{insert}, \texttt{find}, \texttt{erase}, \texttt{save}, \texttt{load}. Быстро стало видно, что основной интерес представляет файловая часть: именно \texttt{save/load} требовали отдельного профилирования. \\
\hline

04.04.2026 &
С помощью \texttt{sample} были сняты профили baseline-версии. Для \texttt{save} большая часть времени наблюдалась внутри \texttt{WriteInOrder} и вызовов \texttt{fwrite}; для \texttt{load} горячими точками оказались \texttt{ReadFromFile}, многочисленные мелкие чтения и последующие вставки в дерево. \\
\hline

04.04.2026 &
При помощи \texttt{leaks} и \texttt{heap} проверена память в сценарии \texttt{load}. Утечек обнаружено не было. Пиковый physical footprint составил около \(34.4\) -- \(34.5\) МБ, а текущий объём занятой памяти во время удержания процесса --- около \(29.6\) -- \(29.7\) МБ. \\
\hline

04.04.2026 &
В \texttt{fixed/main.cpp} реализована внутренняя буферизация чтения и записи: вместо большого числа мелких \texttt{fread/fwrite} добавлены собственные буферизованные классы \texttt{BufferedReader} и \texttt{BufferedWriter}. Формат файла и логика словаря при этом не изменялись. \\
\hline

04.04.2026 &
После исправления повторно выполнены те же профильные замеры и workload-бенчмарки. Получено измеримое ускорение \texttt{save} и \texttt{load}, при этом обычные сценарии с операциями словаря остались на том же уровне производительности. \\
\hline
\end{longtable}

\section{Найденные недочёты}

В ходе исследования был обнаружен один явный прикладной недочёт исходной версии:
\begin{itemize}
\item операции \texttt{save} и \texttt{load} выполняли слишком много мелких обращений к \texttt{fwrite} и \texttt{fread}. Это создавало лишние накладные расходы на файловый ввод-вывод и было хорошо видно по профилям \texttt{sample}.
\end{itemize}

Важно отметить, что ошибок работы с памятью в проверенном сценарии найдено не было:
\begin{itemize}
\item \texttt{leaks} показал \(0\) утечек как для baseline-, так и для fixed-версии;
\item \texttt{heap} показал практически одинаковый объём занятой памяти для обеих версий, то есть исправление не привело к заметному росту потребления памяти.
\end{itemize}

\section{Сравнение версий}

Ниже приведены профильные результаты для сценария из \(200000\) элементов. В таблице сравниваются именно те операции, ради которых вносилось исправление.

\begin{center}
\begin{tabular}{|l|r|r|r|}
\hline
Операция & Baseline, \texttt{us} & Fixed, \texttt{us} & Ускорение \\
\hline
\texttt{save} & \(54708\) & \(28727\) & \(1.90\times\) \\
\hline
\texttt{load} & \(99136\) & \(79050\) & \(1.25\times\) \\
\hline
\end{tabular}
\end{center}

Таким образом, найденный недочёт действительно был существенным: после перехода к внутренней буферизации сохранение стало почти в два раза быстрее, а загрузка ускорилась примерно на четверть. При этом формат файла, внешний интерфейс словаря и логика дерева не менялись, то есть исправление осталось локальным и легко проверяемым.

С точки зрения памяти baseline- и fixed-версии оказались практически одинаковыми: \texttt{leaks} не обнаружил утечек, а по \texttt{heap} обе версии использовали около \(21.7\) МБ пользовательской heap-памяти в рассмотренном сценарии.
\pagebreak
